%=============================================================================
% Alpha_Quantum_Gravity.tex
% Deep Investigation: What the alpha formula tells us about quantum gravity
% Date: 2026-03-23
%=============================================================================
\documentclass[12pt,a4paper]{article}
\usepackage[utf8]{inputenc}
\usepackage{amsmath,amssymb,amsfonts,amsthm}
\usepackage{hyperref}
\usepackage{booktabs}
\usepackage{geometry}
\usepackage{xcolor}
\usepackage{tcolorbox}
\geometry{margin=2.5cm}

\newtheorem{theorem}{Theorem}
\newtheorem{proposition}{Proposition}
\newtheorem{lemma}{Lemma}
\newtheorem{corollary}{Corollary}
\newtheorem{remark}{Remark}
\newtheorem{conjecture}{Conjecture}

\definecolor{dfdgreen}{RGB}{0,128,0}
\definecolor{dfdblue}{RGB}{0,50,180}
\definecolor{dfdred}{RGB}{200,0,0}

\newcommand{\Tr}{\operatorname{Tr}}
\newcommand{\CP}{\mathbb{C}P}
\newcommand{\slp}{\mathfrak{sl}}

\title{What the $\alpha$ Formula Tells Us About Quantum Gravity\\[6pt]
\large A Deep Investigation of Six Interlinked Questions\\[6pt]
\normalsize DFD Research Programme}
\author{Cross-Reference Analysis}
\date{23 March 2026}

\begin{document}
\maketitle
\tableofcontents
\newpage

%=============================================================================
\section{The Formula Under Investigation}
\label{sec:formula}
%=============================================================================

The master formula for the fine-structure constant in DFD:
\begin{equation}
\boxed{\alpha^{-1} = \frac{\pi^{3/2}}{24}\;\Tr(Y^2)\;k_{\max}\,\frac{k_{\max}+3}{k_{\max}+4}\;\left[1 + \frac{7}{80\bigl((k_{\max}+4)^2-1\bigr)}\right]}
\label{eq:master}
\end{equation}
With $\Tr(Y^2) = 10$, $k_{\max} = 60$: $\alpha^{-1} = 137.035\,999\,854$ (residual $5.6 \times 10^{-9}$ from CODATA 2018).

Equivalently, through the exact electroweak formula:
\begin{equation}
\alpha^{-1} = 6\pi\,\kappa_{U(1)}
\label{eq:ew}
\end{equation}
where $\kappa_{U(1)} = 7.26999\ldots$ is computed from the $a_4$ Seeley--DeWitt coefficient of the spectral action on the Toeplitz-truncated internal geometry $\CP^2 \times S^3$.


%=============================================================================
\part{Question 1: Where Are $\hbar$, $G$, and $c$?}
\label{part:q1}
%=============================================================================

\section{Dimensional Analysis of the Formula}

The formula \eqref{eq:master} contains only:
\begin{itemize}
\item $\pi$ (mathematical constant)
\item $\Tr(Y^2) = 10$ (dimensionless SM hypercharge sum)
\item $k_{\max} = 60$ (dimensionless truncation level)
\item Rational numbers ($3, 4, 7, 80, 24$)
\end{itemize}

No $\hbar$, $G$, $c$, $M_P$, or $v$ appear explicitly. This is a \textbf{purely dimensionless} formula yielding a dimensionless number. This is entirely appropriate: $\alpha = e^2/(4\pi\epsilon_0\hbar c)$ is dimensionless, so any formula for it must also be dimensionless.

\section{Where the Fundamental Constants Hide}

\subsection{They are absent by design, not by accident}

The key insight is that \eqref{eq:master} determines $\alpha$ \textbf{without reference to any dimensionful quantity}. This is precisely the claim: $\alpha$ is a topological invariant of the internal space $\CP^2 \times S^3$ combined with the SM particle content.

The formula arises from the spectral action principle:
\begin{equation}
S = \Tr\,f(D^2/\Lambda^2)
\end{equation}
where $D$ is the Dirac operator on the product space $M^4 \times \CP^2 \times S^3$ and $\Lambda$ is the cutoff. The gauge coupling extraction uses the $a_4$ Seeley--DeWitt coefficient, which is a \textbf{dimensionless ratio} of geometric invariants. Concretely:

\begin{equation}
\frac{1}{g_1^2} = \kappa_{U(1)} = \underbrace{\frac{16\pi \cdot (4\pi)^{-7/2} \cdot V_{\rm int}}{12}}_{\pi^{3/2}/24}\;\cdot\; f_0 \;\cdot\; \Tr(Y^2) \;\cdot\; \Lambda^3_{\rm eff}
\end{equation}

\noindent Each factor is dimensionless:
\begin{itemize}
\item $(4\pi)^{-7/2}$: heat-kernel normalisation in 7 internal dimensions
\item $V_{\rm int} = 4\pi^4$: dimensionless volume of $\CP^2 \times S^3$ (in units where the characteristic radius $= 1$)
\item $f_0 = 1$: zeroth moment of the spectral cutoff function
\item $\Lambda^3_{\rm eff} = k_{\max}(k_{\max}+3)/(k_{\max}+4) = 59.0625$: effective mode count from Toeplitz truncation
\end{itemize}

\subsection{Where $\hbar$ and $c$ enter}

The formula for $\alpha$ does not need $\hbar$ or $c$ because $\alpha$ is defined as $e^2/(4\pi\epsilon_0\hbar c)$ --- any theoretical prediction of $\alpha$ must cancel $\hbar$ and $c$ against the dimensions of $e^2$. The spectral action accomplishes this automatically: the heat-kernel expansion is a ratio of traces that yields dimensionless coupling constants.

If one writes out $\alpha$ in terms of measured quantities:
\begin{equation}
\alpha = \frac{e^2}{4\pi\epsilon_0\hbar c}
\end{equation}
this determines $e$ (given $\hbar$, $c$, $\epsilon_0$). DFD's claim is that the \textbf{number} $1/137.036$ is fixed by topology; the \textbf{units} are absorbed into the definition of $e$.

\subsection{Where $G$ and $M_P$ enter}

Newton's constant $G$ enters DFD through a \textbf{separate} relation. The dimensionless invariant (Section 12 of the Unified Theory):
\begin{equation}
\frac{G\hbar H_0^2}{c^5} = \alpha^{57}
\label{eq:alpha57}
\end{equation}
links $G$ to $\alpha$, $\hbar$, $c$, and $H_0$. Equivalently:
\begin{equation}
M_P = \alpha^{-28.5}\;\frac{\hbar H_0}{c^2}
\end{equation}

This means $G$ (or $M_P$) is determined by $\alpha$ plus one cosmological measurement ($H_0$). The $\alpha$ formula itself does \textbf{not} contain $G$ because $\alpha$ is logically prior to $G$ in DFD's axiom structure.

\subsection{The claimed claim: $v = M_P\,\alpha^8\,\sqrt{2\pi}$}

This relation was examined in the Agent 13 analysis (W4i40). Numerical evaluation gives $v \approx 57$ GeV, not the experimental 246 GeV. \textbf{This relation does not hold.} The Higgs VEV is not yet predicted from first principles in DFD.

\begin{tcolorbox}[colback=blue!5!white, colframe=blue!60!black, title=Answer to Question 1]
$\hbar$, $G$, and $c$ do not appear in the $\alpha$ formula because $\alpha$ is dimensionless and its value is determined purely by the topology of $\CP^2 \times S^3$ and the SM particle content. $G$ enters through a \textbf{separate} topological relation $G\hbar H_0^2/c^5 = \alpha^{57}$, making $G$ a derived quantity from $\alpha$ + cosmology. The Higgs VEV $v$ is not yet connected to $M_P$ and $\alpha$ by a verified formula.
\end{tcolorbox}


%=============================================================================
\part{Question 2: The Self-Consistency Equation}
\label{part:q2}
%=============================================================================

\section{Is the Formula Self-Referential?}

\subsection{The claim examined}

The question asks: since $M_P$ enters through $v = M_P\alpha^8\sqrt{2\pi}$, does $\alpha$ depend on itself? The answer is: \textbf{no}, because that relation does not hold (see Part~\ref{part:q1}).

However, a subtler self-consistency question does arise. The $\alpha^{57}$ relation \eqref{eq:alpha57} links $\alpha$, $G$, $\hbar$, $c$, and $H_0$. If we regard $H_0$ as a cosmological input and ask whether the system is consistent, we get:

\begin{equation}
\alpha^{-1} = f(\text{topology, SM content}) = 137.036\ldots
\end{equation}
\begin{equation}
G = \frac{\alpha^{57}\,c^5}{\hbar H_0^2}
\end{equation}

These are \textbf{two independent equations} determining two quantities ($\alpha$ and $G$) from two independent data sources (internal topology and $H_0$). There is no circular self-reference.

\subsection{The formula as a fixed-point equation}

The formula \eqref{eq:master} is \textbf{not} of the form $\alpha = f(\alpha)$. It is a direct evaluation:
\begin{equation}
\alpha^{-1} = \frac{\pi^{3/2}}{24} \times 10 \times 60 \times \frac{63}{64} \times \left[1 + \frac{7}{327600}\right] = 137.036\ldots
\end{equation}

Every input on the right-hand side is either:
\begin{enumerate}
\item A mathematical constant ($\pi$)
\item A topological invariant of $\CP^2 \times S^3$ ($k_{\max} = 60$, dimension $d = 64$)
\item A Standard Model content number ($\Tr(Y^2) = 10$, $N_{\rm sp} = 7$, $g_F = 8$)
\end{enumerate}

None of these depend on $\alpha$ itself. The formula is a \textbf{theorem} of the axiom system A1--A4, not a self-consistency condition.

\subsection{Uniqueness}

Given the axioms, the formula has a unique solution. There is no free parameter to vary; the right-hand side evaluates to a definite number. The only freedom is the \textbf{discrete} choice of Branch A vs.\ Branch B in the microsector fork:
\begin{center}
\begin{tabular}{lcc}
\toprule
\textbf{Branch} & \textbf{Factor} & $\alpha^{-1}$ \\
\midrule
A (regular-module) & $4096/4095$ & 137.03600 \\
B (fermion-rep) & $4095/4096$ & 137.03014 \\
\bottomrule
\end{tabular}
\end{center}
Experiment selects Branch A. The solution is unique once the microsector is fixed.

\subsection{Comparison with other self-consistent coupling predictions}

In asymptotically safe quantum gravity (Eichhorn, Held \& Wetterich, 2018), the fine-structure constant \textbf{is} determined by a fixed-point equation: $\beta_\alpha(\alpha^*, g^*_N) = 0$, where $g^*_N$ is the gravitational fixed-point coupling. This is genuinely self-referential --- $\alpha$ appears on both sides. In that framework, the precise value depends on matter content and computational uncertainties in the gravitational sector.

DFD's approach is structurally different: $\alpha$ is computed from a spectral action trace, not from a fixed point of the RG flow. The spectral action is a one-shot computation, not an iterative fixed-point search.

\begin{tcolorbox}[colback=blue!5!white, colframe=blue!60!black, title=Answer to Question 2]
The $\alpha$ formula is \textbf{not} a self-consistency equation. It is a direct evaluation from topological and SM inputs with no circular dependence on $\alpha$. Given the axioms and the microsector choice (Branch A), the solution is unique: $\alpha^{-1} = 137.03599985$.
\end{tcolorbox}


%=============================================================================
\part{Question 3: The Constraint Field and Planck-Scale Gravity}
\label{part:q3}
%=============================================================================

\section{$\psi$ as Constraint Field: The $A_0$ Analogy}

In QED, the temporal component $A_0$ of the gauge field is a \textbf{constraint}: it has no conjugate momentum, propagates no physical degrees of freedom, and is determined instantaneously by the charge distribution through Gauss's law $\nabla^2 A_0 = -\rho/\epsilon_0$.

In DFD, $\psi$ plays the same role for gravity:
\begin{equation}
\nabla^2\psi = -\frac{4\pi G\rho}{\mu(\chi)\,c^2}
\label{eq:constraint}
\end{equation}
$\psi$ has no propagator, no independent degrees of freedom, and is determined by the matter distribution on each quantum branch.

\section{Why the Spectral Action Determines $\alpha$ Without $\psi$ Propagating}

\begin{proposition}[Constraint Field Decouples from Spectral Action]
\label{prop:decouple}
The spectral action computation for $\alpha^{-1}$ does not require $\psi$ to propagate.
\end{proposition}

\textbf{Argument.} The spectral action computes gauge coupling constants from the $a_4$ Seeley--DeWitt coefficient of the Dirac operator on the \textbf{internal space} $\CP^2 \times S^3$. The constraint field $\psi$ enters only the \textbf{external} 4D conformal factor $e^{2\psi}$ of the metric. Since the heat-kernel expansion on the product space factorises:
\begin{equation}
K(t) = K_{\rm ext}(t) \times K_{\rm int}(t)
\end{equation}
and the gauge coupling extraction uses only $K_{\rm int}(t)$, the constraint field drops out entirely.

$\psi$ is a section of the trivial bundle over $M^4$. The gauge kinetic term arises from the curvature of the internal connection. The two are geometrically independent.

\section{Three Profound Implications for Planck-Scale Gravity}

\subsection{Implication 1: Gravity at the Planck scale is non-propagating}

The fact that $\alpha$ can be computed without $\psi$ propagating means that at the energy scale where the internal geometry determines coupling constants, gravity is \textbf{purely constraint-like}. There is no graviton in DFD. The gravitational field is the conformal factor $e^{2\psi}$, determined by the matter configuration.

This is a radical departure from string theory and loop quantum gravity, both of which posit propagating gravitational degrees of freedom at the Planck scale. In DFD, the Planck scale is not where gravity becomes dynamical --- it is where the internal geometry becomes noncommutative (fuzzy $\CP^2$) and determines the values of all coupling constants.

\subsection{Implication 2: The internal geometry replaces the graviton}

In standard quantum gravity approaches, coupling constant values are inputs. In DFD, they are outputs of the spectral action on the internal space. The role traditionally played by graviton exchange at high energies (providing UV completion) is instead played by the \textbf{spectral geometry of the internal space}:

\begin{center}
\begin{tabular}{lll}
\toprule
\textbf{Feature} & \textbf{Standard QG} & \textbf{DFD} \\
\midrule
Gravity at Planck scale & Propagating graviton & Constraint field \\
Coupling constants & Free parameters & Spectral action outputs \\
UV completion & Graviton loops & Toeplitz truncation \\
Internal space & Compactified dimensions & $\CP^2 \times S^3$ (emergent) \\
Matter-gravity coupling & $\sqrt{G}$ vertex & $\psi$-constraint equation \\
\bottomrule
\end{tabular}
\end{center}

\subsection{Implication 3: The Toeplitz truncation IS the quantum gravity}

The factor $(k_{\max}+3)/(k_{\max}+4) = 63/64$ deviates from unity by $1/64 \approx 1.6\%$. In the continuum limit $k_{\max} \to \infty$, this correction vanishes and $\alpha^{-1} \to (\pi^{3/2}/24) \times \Tr(Y^2) \times k_{\max}$, which diverges. The finite value of $\alpha$ \textbf{requires} the Toeplitz truncation.

This means the finiteness of $\alpha$ is itself a quantum gravity effect: the internal geometry must be quantised (finite matrix algebra $M_{64}(\mathbb{C})$) for the formula to give a finite answer. The truncation at $k_{\max} = 60$ is not a regulator to be removed --- it \textbf{is} the quantum geometry of internal space.

Similarly, the species correction $\varepsilon_w = 1 + 7/(80 \times 4095)$ contains the factor $1/(d^2 - 1) = 1/4095$, which is a finite-dimensional trace effect that vanishes as $d \to \infty$. This is the back-reaction of matter degrees of freedom on the internal geometry --- a quantum gravity effect measurable in the value of $\alpha$.

\subsection{Quantifying the quantum gravity contribution}

We can decompose $\alpha^{-1}$ into its ``classical'' and ``quantum gravity'' parts:
\begin{align}
\alpha^{-1}_{\rm classical} &= \frac{\pi^{3/2}}{24} \times \Tr(Y^2) \times k_{\max} = 139.21\ldots \\
\alpha^{-1}_{\rm QG\,correction\,1} &= -\frac{\pi^{3/2}}{24} \times \Tr(Y^2) \times k_{\max} \times \frac{1}{k_{\max}+4} = -2.18\ldots \\
\alpha^{-1}_{\rm QG\,correction\,2} &= +\frac{\pi^{3/2}}{24} \times \Tr(Y^2) \times k_{\max} \times \frac{k_{\max}+3}{k_{\max}+4} \times \frac{7}{80(d^2-1)} = +0.003\ldots
\end{align}

The first quantum gravity correction (finite-size tracelessness) shifts $\alpha^{-1}$ by $-2.18$, a $1.6\%$ effect. The second correction (microsector trace) shifts by $+0.003$, a $2 \times 10^{-5}$ effect. Both are finite, calculable, and essential for matching experiment.

\begin{tcolorbox}[colback=blue!5!white, colframe=blue!60!black, title=Answer to Question 3]
Since $\psi$ is a constraint field, the spectral action on $\CP^2 \times S^3$ determines $\alpha$ \textbf{without} $\psi$ propagating. This means:
\begin{enumerate}
\item Gravity at the Planck scale is non-propagating (no graviton).
\item The internal geometry replaces the graviton as the UV completion mechanism.
\item The Toeplitz truncation ($k_{\max} = 60$) IS the quantum gravity --- without it, $\alpha$ would diverge. The finiteness of $\alpha$ is literally a quantum gravity prediction.
\item The quantum gravity correction to $\alpha^{-1}$ is $-2.18$ (1.6\%), coming from the finite-size tracelessness factor $63/64$.
\end{enumerate}
\end{tcolorbox}


%=============================================================================
\part{Question 4: BMV Prediction and the $\alpha$ Formula}
\label{part:q4}
%=============================================================================

\section{The BMV Enhancement Prediction}

The Bose--Marletto--Vedral (BMV) experiment proposes to create gravitational entanglement between two mesoscopic masses in spatial superposition. DFD predicts a $\sim 2200\times$ enhancement of the gravitational entanglement rate compared to the standard Newtonian prediction, arising from the MOND interpolation function $\mu(\chi)$ in the low-acceleration regime.

The enhancement factor scales as $1/\mu(\chi)$ where $\chi = |\nabla\psi|^2/a_*^2$ is the dimensionless acceleration parameter. For the BMV experimental regime where accelerations are $a \ll a_0$:
\begin{equation}
\text{Enhancement} = \frac{1}{\mu(\chi)} \approx \frac{a_0}{a} \approx 2200
\end{equation}

\section{Does the $\alpha$ Formula Constrain the BMV Prediction?}

\subsection{Direct connection through $\Tr(Y^2)$}

The factor $\Tr(Y^2) = 10$ appears in both:
\begin{enumerate}
\item The $\alpha$ formula: $\alpha^{-1} \propto \Tr(Y^2)$
\item The MOND scale: $a_0 = 2\sqrt{\alpha}\,cH_0$ (derived relation)
\end{enumerate}

The connection is \textbf{indirect but significant}. The $\alpha$ formula determines $\alpha$ from the SM spectrum. The MOND scale $a_0$ depends on $\alpha$ through the derived relation. And the BMV enhancement depends on $a_0$.

The chain is:
\begin{equation}
\text{SM spectrum} \;\xrightarrow{\text{spectral action}}\; \alpha \;\xrightarrow{a_0 = 2\sqrt{\alpha}\,cH_0}\; a_0 \;\xrightarrow{1/\mu}\; \text{BMV enhancement}
\end{equation}

\subsection{Quantitative constraint}

The precision of the $\alpha$ formula ($5.6 \times 10^{-9}$) constrains BSM contributions to $\Tr(Y^2)$ at the level:
\begin{equation}
|\Delta\Tr(Y^2)| < 5.6 \times 10^{-8}
\end{equation}

Any new light species that modified $\Tr(Y^2)$ would change both $\alpha$ and (through $a_0$) the MOND scale. The $\alpha$ formula therefore provides an extraordinarily tight constraint on the matter content that enters the BMV prediction.

\subsection{Does the enhancement factor depend on $\alpha$ directly?}

The MOND interpolation function $\mu(\chi)$ is determined by the DFD Lagrangian $\mathcal{L} = \Lambda_\psi^4[\chi - \ln(1+\chi)]$, which depends on $a_* = cH_0/\sqrt{\alpha \times k_a}$ with $k_a = 3/(8\alpha)$. Therefore:
\begin{equation}
a_* = cH_0\sqrt{\frac{8}{3}} \propto cH_0
\end{equation}
The $\alpha$ dependence cancels in $a_*$ (since $k_a \propto 1/\alpha$), and $a_0 = 2\sqrt{\alpha}\,cH_0$ carries the residual dependence. The BMV enhancement $\sim a_0/a \propto \sqrt{\alpha}$ depends weakly on $\alpha$: a 1\% change in $\alpha$ would shift the enhancement by $0.5\%$. Given the sub-ppm precision of the $\alpha$ formula, this contribution to the BMV uncertainty is negligible compared to experimental uncertainties (mass, coherence time, acceleration environment).

\begin{tcolorbox}[colback=blue!5!white, colframe=blue!60!black, title=Answer to Question 4]
The $\alpha$ formula constrains the BMV prediction through a chain: $\Tr(Y^2) \to \alpha \to a_0 \to$ enhancement. The constraint is powerful: any BSM physics altering $\Tr(Y^2)$ by more than $6 \times 10^{-8}$ would violate the $\alpha$ match. The direct dependence of the $2200\times$ enhancement on $\alpha$ is weak ($\propto \sqrt{\alpha}$), so the sub-ppm $\alpha$ precision provides negligible uncertainty compared to experimental factors. The main role of the $\alpha$ formula is to \textbf{lock the SM content} that enters the MOND-scale prediction.
\end{tcolorbox}


%=============================================================================
\part{Question 5: The Born Rule and High-Energy Modifications}
\label{part:q5}
%=============================================================================

\section{The DFD Born Rule Derivation}

DFD derives the Born rule $p_i = |\alpha_i|^2$ from three ingredients:

\begin{enumerate}
\item \textbf{Gravitational einselection (i34):} The constraint equation $\nabla^2\psi = -4\pi G\rho/(\mu c^2)$ selects the position basis as the preferred basis. Different mass configurations source orthogonal $\psi$ fields: $\langle[\psi]_i|[\psi]_j\rangle = 0$ for $\rho_i \neq \rho_j$.

\item \textbf{Envariance (after Zurek):} For entangled states $|\Psi\rangle = \sum_i \alpha_i |s_i\rangle|e_i\rangle|[\psi]_i\rangle$, the probability assignment $p_i = |\alpha_i|^2$ is the unique measure compatible with envariant symmetry + additivity.

\item \textbf{Gleason's theorem:} The constraint field enlarges the effective Hilbert space dimension beyond 2, closing the qubit loophole.
\end{enumerate}

The probability density in the physical (conformal) frame is:
\begin{equation}
\rho_{\rm prob} = e^{2\psi}|\Psi|^2
\label{eq:born-dfd}
\end{equation}
where the factor $e^{2\psi}$ arises from the conformal metric $\tilde{g}_{\mu\nu} = e^{2\psi}g_{\mu\nu}$ (the physical metric determinant contributes $\sqrt{-\tilde{g}} = e^{4\psi}\sqrt{-g}$ in 4D, giving $e^{2\psi}$ per spatial dimension in the probability measure).

\section{Does the $\alpha$ Formula Modify the Born Rule at High Energies?}

\subsection{Structural independence}

The $\alpha$ formula and the Born rule derivation operate on \textbf{different geometric sectors}:

\begin{center}
\begin{tabular}{lll}
\toprule
& \textbf{Born rule} & \textbf{$\alpha$ formula} \\
\midrule
Operates on & External $M^4$ & Internal $\CP^2 \times S^3$ \\
Uses & Constraint eq.\ $\nabla^2\psi = \ldots$ & Spectral action $a_4$ \\
Field & $\psi$ (external conformal) & Internal connection \\
Depends on $\alpha$ & No & Determines $\alpha$ \\
\bottomrule
\end{tabular}
\end{center}

The Born rule derivation uses the \textbf{external} constraint equation, while the $\alpha$ formula uses the \textbf{internal} spectral geometry. These are geometrically factored on the product space $M^4 \times \CP^2 \times S^3$, so modifications to one do not directly affect the other.

\subsection{Could high-energy corrections modify the Born rule?}

At energies approaching the Toeplitz cutoff scale $\Lambda_{\rm Toeplitz} \sim k_{\max}^{1/2} M_P$, two effects could modify the Born rule:

\begin{enumerate}
\item \textbf{$\psi$-dependent vacuum polarization:} At sufficiently high energies, vacuum polarization loops involving the constraint field could modify the effective probability measure. However, these contributions are suppressed by $\kappa^2 \sim (M_P/E)^4$ at sub-Planckian energies and are unobservable in any foreseeable experiment.

\item \textbf{Noncommutative corrections from the internal space:} If the Toeplitz truncation affects the factorisation of the heat kernel at very high energies, cross-terms between internal and external sectors could arise. These would be $O(1/k_{\max}^2) \sim 3 \times 10^{-4}$ corrections to the Born rule at the Planck scale.
\end{enumerate}

\subsection{The critical distinction: $e^{2\psi}$ vs.\ $|\Psi|^2$}

The DFD Born rule \eqref{eq:born-dfd} has the structure $\rho = e^{2\psi}|\Psi|^2$. The $e^{2\psi}$ factor is a \textbf{conformal volume element}, not a modification of quantum probabilities. In a region with $\psi \neq 0$, the physical volume element differs from the coordinate volume element. The quantum probability $|\Psi|^2$ is unchanged; only the measure on which it is integrated changes.

This distinction is crucial: the $\alpha$ formula does not modify $|\Psi|^2$ because $\alpha$ is determined by the internal geometry while $|\Psi|^2$ is determined by the quantum state on the external space.

\subsection{The vector extension caveat}

The i39 analysis (W4i39\_ThirdPeak\_Resolution.tex) identified a potential issue: if DFD is extended to include a vector field component (for CMB compatibility), the constraint algebra changes, and the Born rule derivation from the scalar $\psi$ constraint field ``may be modified.'' This is rated \textcolor{dfdred}{\textbf{YELLOW}} --- the Born rule derivation must be re-examined in any vector extension. However, this is a concern about extending DFD beyond its current axioms, not about the $\alpha$ formula modifying the existing Born rule.

\begin{tcolorbox}[colback=blue!5!white, colframe=blue!60!black, title=Answer to Question 5]
The $\alpha$ formula does \textbf{not} modify the Born rule at any energy. The two operate on geometrically factored sectors (internal vs.\ external). The $e^{2\psi}$ in the DFD probability density is a conformal volume element, not a modification of quantum probabilities. At extremely high energies ($E \sim M_P$), $O(1/k_{\max}^2) \sim 10^{-4}$ corrections from cross-terms between internal and external heat kernels are conceivable but unobservable. The main theoretical risk to the Born rule comes from potential vector-field extensions of DFD, not from the $\alpha$ formula.
\end{tcolorbox}


%=============================================================================
\part{Question 6: Literature Connections}
\label{part:q6}
%=============================================================================

\section{Competing Approaches to Computing $\alpha$ from Quantum Gravity}

\subsection{Asymptotic safety (Eichhorn, Held \& Wetterich, 2018)}

In asymptotically safe quantum gravity, the fine-structure constant is determined by a UV fixed point of the combined gravity-matter RG flow. The beta function for the gauge coupling receives gravitational contributions:
\begin{equation}
\beta_\alpha = \beta_\alpha^{\rm matter} + f_g \cdot g_N \cdot \alpha
\end{equation}
where $g_N$ is the dimensionless Newton coupling and $f_g$ is a gravitational fluctuation coefficient. At the UV fixed point $\beta_\alpha = 0$, the gauge coupling is determined:
\begin{equation}
\alpha^* = -\frac{\beta_\alpha^{\rm matter}}{f_g \cdot g_N^*}
\end{equation}

\textbf{Key difference from DFD:} In asymptotic safety, $\alpha$ is genuinely self-referential (it appears on both sides of the fixed-point equation). The precise value depends on the gravitational fluctuation coefficient $f_g$, which has significant computational uncertainties. In DFD, $\alpha$ is computed from a one-shot spectral action trace with no self-reference.

\textbf{Status (2018):} The asymptotic safety prediction is consistent with $\alpha \sim 1/137$ for SM matter content but lacks the precision to distinguish the correct value at the ppm level. DFD achieves $5.6 \times 10^{-9}$ relative precision.

\subsection{Noncommutative geometry (Chamseddine \& Connes, 1997--2025)}

DFD's spectral action computation is built directly on the Chamseddine--Connes framework. The key results from NCG:

\begin{enumerate}
\item \textbf{$\sin^2\theta_W = 3/8$ at unification} (Chamseddine \& Connes, 1997): the Weinberg angle is predicted from the spectral triple.
\item \textbf{Higgs mass prediction}: $m_H \approx 170$ GeV at tree level, reduced to $\sim 125$ GeV with $\sigma$-field extension.
\item \textbf{Gauge coupling unification}: the spectral action gives the correct GUT structure.
\end{enumerate}

\textbf{What DFD adds:} Standard NCG gives coupling constants at the \textbf{unification scale}. DFD extends this to the \textbf{low-energy} electromagnetic coupling by:
\begin{itemize}
\item Choosing the specific internal space $\CP^2 \times S^3$ (7-dimensional)
\item Implementing Toeplitz truncation at $k_{\max} = 60$
\item Computing the $a_4$ coefficient with all finite-size corrections
\item Resolving the microsector fork (Branch A: regular module)
\end{itemize}

The precision improvement over standard NCG is enormous: from $\sim 10\%$ (GUT-scale coupling ratios) to $5.6 \times 10^{-9}$ (low-energy $\alpha$).

\subsection{String theory landscape}

In string theory, the fine-structure constant is determined by the vacuum expectation value of the dilaton and the topology of the compactification manifold. The string landscape contains $\sim 10^{500}$ vacua, with $\alpha$ varying across the landscape. This is often combined with anthropic selection to ``explain'' the observed value.

\textbf{Key difference from DFD:} String theory has a vast landscape of solutions, requiring anthropic selection. DFD has a \textbf{unique} prediction (once the internal space and SM content are fixed). There is no landscape; the answer is 137.036 or the theory is wrong.

\subsection{Other approaches}

\begin{itemize}
\item \textbf{Loop quantum gravity:} Does not currently predict coupling constant values. The focus is on quantum geometry (area/volume spectra) rather than particle physics parameters.

\item \textbf{Causal dynamical triangulations (CDT):} Provides a non-perturbative definition of the gravitational path integral but has not yet been extended to include matter couplings at the level needed to predict $\alpha$.

\item \textbf{Bivector Standard Model (2025):} A geometric algebra approach where $\alpha$ emerges as a structural ratio related to the electron's internal geometry. Achieves qualitative agreement but not the ppm-level precision of DFD.
\end{itemize}

\section{The Unique Feature of DFD's Prediction}

DFD's $\alpha$ prediction is distinguished from all competitors by three features simultaneously:

\begin{enumerate}
\item \textbf{Precision:} $5.6 \times 10^{-9}$ relative accuracy --- the highest-precision prediction of $\alpha$ from any quantum gravity framework.
\item \textbf{No free parameters:} Zero continuous parameters. The only freedom is the discrete microsector choice, which is falsified for Branch B.
\item \textbf{No self-reference:} The formula is a direct evaluation, not a fixed-point or consistency equation.
\end{enumerate}

No other approach to quantum gravity achieves all three simultaneously.

\begin{tcolorbox}[colback=blue!5!white, colframe=blue!60!black, title=Answer to Question 6]
The main competing approaches are:
\begin{itemize}
\item \textbf{Asymptotic safety}: genuinely self-referential, $\sim 10\%$ precision, significant computational uncertainties.
\item \textbf{Standard NCG}: gives GUT-scale ratios ($\sin^2\theta_W = 3/8$) but not low-energy $\alpha$ directly.
\item \textbf{String landscape}: $\sim 10^{500}$ vacua, requires anthropic selection, no unique prediction.
\item \textbf{LQG/CDT}: not yet extended to predict coupling constants.
\end{itemize}
DFD is unique in achieving $5.6 \times 10^{-9}$ precision with zero free parameters and no self-reference.
\end{tcolorbox}


%=============================================================================
\part{Synthesis: What the $\alpha$ Formula Reveals About Quantum Gravity}
\label{part:synthesis}
%=============================================================================

\section{The Big Picture}

The $\alpha$ formula tells us something profound about the nature of quantum gravity in DFD. We can summarize it in five statements:

\begin{enumerate}
\item \textbf{Coupling constants are spectral invariants of the internal geometry.} The value $\alpha^{-1} = 137.036$ is not a free parameter but a geometric property of $\CP^2 \times S^3$, analogous to how $\pi$ is a geometric property of a circle.

\item \textbf{Quantum gravity means finite matrix geometry, not quantized gravitons.} The Toeplitz truncation at $k_{\max} = 60$ replaces the continuum internal space with a finite matrix algebra $M_{64}(\mathbb{C})$. This is the quantum gravity of DFD: spacetime becomes noncommutative (fuzzy) at the Planck scale, with $d^2 - 1 = 4095$ internal degrees of freedom.

\item \textbf{The finiteness of $\alpha$ is itself a quantum gravity prediction.} Without the finite-size correction $(k_{\max}+3)/(k_{\max}+4)$, the formula gives a divergent result ($\alpha^{-1} \to \infty$ as $k_{\max} \to \infty$). The finite value of $\alpha$ requires the internal geometry to be quantized.

\item \textbf{Gravity and gauge forces have a common origin.} Both emerge from the spectral action on the product space $M^4 \times \CP^2 \times S^3$. The gravitational sector ($\psi$ as constraint on $M^4$) and the gauge sector (connections on $\CP^2 \times S^3$) are geometrically factored but share the same underlying spectral geometry.

\item \textbf{The constraint-field nature of $\psi$ is essential.} If $\psi$ were dynamical (propagating graviton), it would contribute to vacuum polarization loops and modify $\alpha$ through running. As a constraint field, it decouples from the spectral action entirely, allowing $\alpha$ to be determined purely by the internal geometry. The non-propagation of gravity is not a bug --- it is what makes the prediction possible.
\end{enumerate}

\section{Open Questions}

\begin{enumerate}
\item \textbf{Origin of $k_{\max} = 60$:} The Bridge Lemma gives $k_{\max} = \chi(\CP^2, \mathcal{O}(9) \oplus \mathcal{O}^{\oplus 5}) = 60$, and the finite-symmetry closure (icosahedral $A_5$) provides an alternative route. But can $k_{\max}$ be derived from a completeness argument on the Spin$^c$ Dirac operator alone?

\item \textbf{Why Branch A (regular module)?} The microsector choice $\mathcal{H}_F = M_d(\mathbb{C})$ vs.\ $\mathcal{H}_F = \mathbb{C}^d$ is experimentally decided by the sub-ppm match. But is there a first-principles argument that selects the regular module?

\item \textbf{Extension to $\alpha_s$ and fermion masses:} The spectral action gives couplings at the cutoff scale. Can the running be incorporated spectrally (truncation level as RG scale)?

\item \textbf{Connection to the Born rule at the Planck scale:} At energies where the Toeplitz truncation becomes relevant, do cross-terms between internal and external heat kernels modify the probability measure?

\item \textbf{The $v = M_P \alpha^n \times (\ldots)$ relation:} Is there a topological formula connecting the Higgs VEV to $M_P$ and $\alpha$? The $\alpha^8$ attempt fails numerically. What is the correct exponent, if any?
\end{enumerate}

\section{Testable Consequences}

\begin{center}
\renewcommand{\arraystretch}{1.3}
\begin{tabular}{llll}
\toprule
\textbf{Prediction} & \textbf{Source} & \textbf{Test} & \textbf{Timeline} \\
\midrule
$\alpha^{-1} = 137.03599985$ & Spectral action & Cs/Rb atom interferometry & Current \\
$H_0 = 72.09$ km/s/Mpc & $\alpha^{57}$ relation & JWST distance ladder & 2026--2028 \\
No graviton (constraint field) & A1--A2 axioms & BMV null result & 2027--2030 \\
$2200\times$ BMV enhancement & MOND regime & BMV with isolation & 2030+ \\
$\Sigma m_\nu = 61.4$ meV & Spectral action & DESI DR3 & 2027 \\
Zero scalar dipole radiation & Constraint field & SKA pulsar timing & 2028+ \\
\bottomrule
\end{tabular}
\end{center}


%=============================================================================
\section*{References}
%=============================================================================

\begin{enumerate}
\item[{[1]}] G.~Alcock, ``Ab Initio Derivation of the Fine Structure Constant from Density Field Dynamics'' (2025). \url{https://doi.org/10.5281/zenodo.19173548}

\item[{[2]}] G.~Alcock, ``Density Field Dynamics: A Complete Unified Theory,'' v108, Zenodo (2026). Sections 8C, 12, 13; Appendices F, K.

\item[{[3]}] A.~Eichhorn, A.~Held \& C.~Wetterich, ``Quantum-gravity predictions for the fine-structure constant,'' Phys.\ Lett.\ B \textbf{782}, 198 (2018). arXiv:1711.02949.

\item[{[4]}] A.~H.~Chamseddine \& A.~Connes, ``The Spectral Action Principle,'' Comm.\ Math.\ Phys.\ \textbf{186}, 731 (1997).

\item[{[5]}] A.~Connes, ``Noncommutative Geometry, the Spectral Standpoint,'' arXiv:1910.10407 (2019).

\item[{[6]}] S.~Bose, A.~Mazumdar et al., ``Spin entanglement witness for quantum gravity,'' PRL \textbf{119}, 240401 (2017).

\item[{[7]}] C.~Marletto \& V.~Vedral, ``Gravitationally induced entanglement between two massive particles is sufficient evidence of quantum effects of gravity,'' PRL \textbf{119}, 240402 (2017).

\item[{[8]}] D.~L.~Danielson, G.~Satishchandran \& R.~M.~Wald, ``Gravitationally mediated entanglement: Newtonian field vs.\ graviton,'' PRD \textbf{105}, 086001 (2022).

\item[{[9]}] W.~H.~Zurek, ``Quantum Theory of the Classical: Einselection, Envariance, Quantum Darwinism and Extantons,'' Entropy \textbf{24}, 1520 (2022).

\item[{[10]}] A.~H.~Chamseddine, A.~Connes \& W.~D.~van Suijlekom, ``Unified Pati--Salam from Noncommutative Geometry,'' arXiv:2511.07672 (2025).

\item[{[11]}] J.~Barrett, ``Hearing the Shape of the Universe,'' arXiv:2511.05909 (2025).
\end{enumerate}


\end{document}
